\begin{proposition}[Every torus-degeneration fibre has even Euler number]\hypertarget{tag-01CM}{}\label{cdp:st:C.2.UN.1}
\StatementRecord{cdp:st:C.2.UN.1}{proved}{01CM}{\SeriesRef{cdp:eq:C.UN.3}, \SeriesRef{cdp:eq:C.UN.4}}

Let $f:Y\to B$ be a proper holomorphic map from a smooth complex
threefold to a smooth complex curve, with connected fibres and smooth
general fibre a complex two-torus. For every $b\in B$,
\[
                    e((Y_b)_{\rm red})\in2\Z.       \tag{C.UN.1}\hypertarget{tag-0152}{}\label{cdp:eq:C.UN.1}
\]
The fibres may have arbitrary multiplicities and singularities.
\end{proposition}


\begin{proof}

\begin{enumerate}[label=\arabic*.,leftmargin=*]
\item Take a sufficiently small closed coordinate disc $\Delta$ about $b$,
containing no other critical value. Put $U=f^{-1}(\Delta)$ and
$M=\partial U$. Then $U$ is a compact oriented smooth real six-manifold
with boundary, and $M$ is the mapping torus of the orientation-preserving
monodromy of a real four-torus. The analytic central support $S$ is a
deformation retract of $U$. One precise construction uses a controlled
gradient for the real analytic function $|f|^2$ on a Whitney
stratification: it decreases $|f|$, is tangent to the strata over zero,
and has trajectories with a limit on the zero set by the local
analytic gradient inequality. Compactness of the zero set gives a
common small disc, and the resulting continuous limit homotopy fixes
the zero set. Equivalently the triangulated proper local tube is a
regular neighbourhood of the zero fibre. Nilpotents do not change this
retraction. Thus $e(S)=\chi(U)$.

\item We give the parity calculation entirely over $\Q$. Write
$s_i=\operatorname{rank}(H_i(U)\to H_i(U,M))$.
Exactness of the pair sequence, with Poincar\'e--Lefschetz duality,
gives
\[
 b_0(M)+b_1(M)+b_2(M)
       =\sum_{i=0}^{5}b_i(U)-s_0-s_3-2s_1-2s_2.
                                                        \tag{C.UN.2}\hypertarget{tag-0153}{}\label{cdp:eq:C.UN.2}
\]
Indeed, the summand $b_i(M)$ is the sum of the ranks of
$H_i(M)\to H_i(U)$ and $H_{i+1}(U,M)\to H_i(M)$;
substitution of $b_{i+1}(U,M)=b_{5-i}(U)$ gives the formula.
Here $s_0=0$ and $b_6(U)=0$. The map defining $s_3$ is represented,
after duality, by the middle-dimensional intersection pairing of $U$.
That pairing is skew-symmetric because the middle dimension is three,
so $s_3$ is even. Consequently
\[
 \chi(U)\equiv b_0(M)+b_1(M)+b_2(M)\pmod2.            \tag{C.UN.3}\hypertarget{tag-0154}{}\label{cdp:eq:C.UN.3}
\]

\item Let $T$ act on $H^1(F,\Q)$ for the four-torus fibre $F$, and set
$a_i=\dim\ker(\Lambda^iT-I)$. Cutting the mapping torus along one
fibre gives the Wang exact sequence; hence
\[
 b_0(M)=1,\quad b_1(M)=1+a_1,\quad b_2(M)=a_1+a_2.
                                                        \tag{C.UN.4}\hypertarget{tag-0155}{}\label{cdp:eq:C.UN.4}
\]
Orientation preservation says $\det T=1$. The transformation
$G=\Lambda^2T$ preserves the nondegenerate symmetric wedge pairing
on $\Lambda^2\Q^4$ and has determinant $(\det T)^3=1$.
The rank $\operatorname{rank}(G-I)$ is even. Complexify and split off the generalized
$(-1)$-eigenspace. Its dimension is even: all eigenvalues outside
$\{1,-1\}$ occur in reciprocal pairs with equal algebraic
multiplicity, and $\det G=1$. This generalized eigenspace is
nondegenerate and orthogonal to its invariant complement. On it
$G-I$ is invertible, of even rank. On the complement $G+I$ is
invertible, and
\[
                         (G-I)(G+I)^{-1}
\]
is skew-adjoint for the symmetric wedge pairing. Multiplication by
that nondegenerate pairing represents it as a skew-symmetric matrix,
whose rank is even. Its rank equals that of $G-I$ on the complement.
Since $\dim\Lambda^2\Q^4=6$, $a_2$ is
even. Equations \eqref{cdp:eq:C.UN.3}--\eqref{cdp:eq:C.UN.4} prove \eqref{cdp:eq:C.UN.1}.

\end{enumerate}
\end{proof}